
% ============================
% CHAPTERS 6--10 FOR THE TOME
% ============================

\part{Core Field Dynamics and Gravitational Structure}

\chapter{ER Network Formalism}

\section{Planck-Scale Adjacency Graph}

We model microscopic spatial adjacency not by metric distance but by 
an Einstein--Rosen (ER) connectivity graph:
\[
W_{ij} \ge 0,\qquad W_{ij}=W_{ji},
\]
where $i,j$ label Planck-scale sites.  
Edges represent wormhole-like entanglement connections.

\section{ER Contribution to Geometry}

We define an effective ER metric component:
\[
g^{(ER)}_{\mu\nu} = \lambda 
\sum_{i,j} W_{ij}\, \xi^{(i)}_\mu \xi^{(j)}_\nu,
\]
where $\xi^{(i)}_\mu$ projects fields onto site neighborhoods.

\section{Entanglement and Curvature}

The Bianchi-compatibility condition suggests:
\[
R_{\mu\nu} \sim \nabla_\mu\nabla_\nu S_{\mathrm{ER}},
\]
where $S_{\mathrm{ER}}$ is ER-network entanglement entropy.


\chapter{Phase--Chromatic Coupling}

\section{Interaction Term}

We introduce a coupling Lagrangian:
\[
\mathcal{L}_{\tau\chi}
=
-\gamma\,
G^{ab}(\bm{\chi})\, 
(\partial_\mu \tau)(\partial^\mu \chi_a)\chi_b .
\]

\section{Coupled Equations}

Phase equation acquires a chromatic backreaction:
\[
\beta \Box \tau - V'(\tau)
=
\gamma\, G^{ab}\chi_b\,\Box\chi_a .
\]

Chromatic equations acquire phase modulation:
\[
\nabla_\mu(G^{ab}\partial^\mu\chi_b)
- \partial^a U
=
\gamma\,(\partial_\mu\tau)(\partial^\mu\chi_a).
\]


\chapter{Planck-Scale Tunneling and Coherence}

\section{Tunneling Regime}

As $\lambda_\gamma\to \ell_P$, photon modes lose separability.  
We describe tunneling amplitude:
\[
\mathcal{T} 
\sim 
\exp\left(-\frac{\Delta S}{\hbar}\right),
\]
where $\Delta S$ is the action separation between adjacent bead states.

\section{Coherence Condition}

Bead-lattice coherence requires:
\[
|\partial_\mu\tau| < \tau_c^{-1},
\]
ensuring phase gradients do not decohere chromatic coupling.

\section{Gravity from Tunneling}

In the coherent limit:
\[
G_{\mu\nu}^{(\text{eff})}
\propto
\langle \partial_\mu \tau \partial_\nu\tau \rangle_{\rm tunnel}.
\]


\chapter{White-State Core Dynamics}

\section{White-State Definition}

The white-state core $\Wcore$ satisfies:
\[
\bm{\chi}=0,\qquad 
\nabla_\mu \tau =0,\qquad 
W_{ij}=\delta_{ij}.
\]

\section{Stability Condition}

Linearizing perturbations:
\[
\delta \bm{\chi}=0,\qquad 
\partial_\mu \delta\tau =0
\]
remains invariant under the evolution operator.

\section{Core Geometry}

Effective curvature:
\[
R_{\mu\nu}(\Wcore)
=
\Lambda_{\rm eff}\, g_{\mu\nu},
\]
with $\Lambda_{\rm eff}$ determined by chromatic vacuum energy.


\chapter{Collapse, Bounce, and Nonsingular Evolution}

\section{Collapse to $\Wcore$}

Under extreme compression:
\[
\bm{\chi}\to 0,
\qquad
\partial_\mu\tau \to 0.
\]

\section{Bounce Condition}

The bounce is triggered when:
\[
\rho_{\chi} + \rho_\tau + \rho_{ER}
=
\rho_{\rm crit},
\]
with $\rho_{\rm crit}$ finite.

\section{Re-expansion}

A new ER foam emerges:
\[
W_{ij}^{(\text{new})} \sim W_{ij}^{(\text{old})} + \epsilon_{ij},
\]
seeding the next cosmic cycle.
